\input{../tex-inputs/latex-leseansicht-vorspann}

\section[Gustav Schwarzkopf an Arthur Schnitzler, 5. 6. 1913]{L04256 Gustav Schwarzkopf an Arthur Schnitzler, 5. 6. 1913}
\nopagebreak\mylabel{L04256v}
\rehead{ }\normalsize\beginnumbering\briefempfaengerindex{Schnitzler, Arthur@\textsc{Schnitzler, Arthur}!zzzSchwarzkopf, Gustav@\emph{von Gustav Schwarzkopf}!1913-06-051@{5. 6. 1913}|(be}
\toendnotes[C]{\smallbreak\pagebreak[2]}
\correspDesc{Versand  durch Gustav Schwarzkopf am 5. 6. 1913 in Baden bei Wien
\newline{}Erhalt  durch Arthur Schnitzler im Zeitraum [6. 6. 1913 – 10. 6. 1913?] in Wien}\toendnotes[C]{\smallbreak}
\Standort{CUL, Schnitzler, B 96a.}
\physDesc{Bildpostkarte, 405 Zeichen
\newline{}Handschrift: schwarze Tinte, deutsche Kurrent
\newline{}Versand: Stempel: »\nobreak{}\oindex{Baden bei Wien@\textbf{Baden bei Wien}, \emph{Hauptstadt}|pwk}Baden, N. Ö. 4, {[}5.{]} VI. {[}1{]}3\nobreak{}«.  }\toendnotes[C]{\smallbreak}\pstart{}{\pb}Herrn\pend{}\pstart{}D\textsuperscript{r} Arthur Schnitzler\pend{}\pstart{}\textsc{Wien XVIII.\oindex{XVIII., Währing@\textbf{XVIII., Währing}, \emph{Verwaltungsgebiet}|pw}}\pend{}\pstart{}Sternwarteſtraße 71\oindex{Wien@\textbf{Wien}!XVIII., Währing@\textbf{XVIII., Währing}!Sternwartestraße 71@\textbf{Sternwartestraße 71}, \emph{Wohngebäude}|pw}\pend{}{\bigskip}
\pstart
           \noindent{}\centering{}{\pb}\textcolor{gray}{\textbf{BADEN – Helenenthal, Ruine Rauhenstein}}\oindex{Burgruine Rauhenstein@\textbf{Burgruine Rauhenstein}, \emph{Ruine}|pw}\pend
           \vspace{1em}
\pstart
           \raggedleft{}{\pb}5. VI. 13.\pend
           
\pstart{}Lieber Arthur!\pend\vspace{0.5em}
\pstart
           Es iſt zu heiß um nach Wien\oindex{Wien@\textbf{Wien}, \emph{Verwaltungsgebiet}|pw} zu fahren, ich will
               noch einen Tag abwarten. Ihrer freundliche Einladung für \label{K_L04256-1v}\edtext{Freitag{ }\uline{Abend}}{\lemma{\textnormal{\emph{Freitag Abend}}}\Cendnote{\textnormal{Vgl. A. S.: \emph{Tagebuch}, 30. 5. 1913.}}}\label{K_L04256-1} ko{\geminationn}te ich doch
               nicht folgen; gesellschafts\uline{fähig} – ſo für einen
               größeren Kreis – bin ich doch noch nicht. \label{K_L04256-2v}\edtext{So{\geminationn}tag bin ich
               gewiß{ }ſchon in Wien\oindex{Wien@\textbf{Wien}, \emph{Verwaltungsgebiet}|pw} und{ }ſuche Sie dann Nachmittag}{\lemma{\textnormal{\emph{Sonntag … Nachmittag}}}\Cendnote{\textnormal{Es wurde der Abend, vgl. A. S.: \emph{Tagebuch}, 8. 6. 1913.}}}\label{K_L04256-2}
               auf.\pend
           
\pstart
           Ihre Damen\pwindex{Schnitzler, Olga 17.\,1.\,1882 Wien – 13.\,1.\,1970 Lugano@\textsc{Schnitzler, Olga} (17.\,1.\,1882 Wien – 13.\,1.\,1970 Lugano), \emph{Schauspielerin, Sängerin}|pwv}\pwindex{Cappellini, Lili 13.\,9.\,1909 Wien – 26.\,7.\,1928 Venedig@\textsc{Cappellini, Lili} (13.\,9.\,1909 Wien – 26.\,7.\,1928 Venedig)|pwv} und{\\[\baselineskip]}Sie
               herzlichſt grüßend{\\[\baselineskip]}Ihr{\\[\baselineskip]}Gustav\pend
           \leftskip=0em{}\selectlanguage{ngerman}\endnumbering\briefempfaengerindex{Schnitzler, Arthur@\textsc{Schnitzler, Arthur}!zzzSchwarzkopf, Gustav@\emph{von Gustav Schwarzkopf}!1913-06-051@{5. 6. 1913}|)be}\mylabel{L04256h}
\begin{anhang}
\end{anhang}\newcommand{\dateiname}{L04256}\newcommand{\titel}{Gustav Schwarzkopf an Arthur Schnitzler, 5. 6. 1913}\newcommand{\editorInnen}{Herausgegeben von Jahnke, SelmaMüller, Martin Anton}\input{../tex-inputs/latex-leseansicht-abspann}
